%%%%%%%%%%%%%%%%%%%%%%%%%%%%%%%%%%%%%%%%%%%%%%%%%%%%%%%%%%%%%%%%%%%%%%%%%
%
% Plantilla para libro de texto de matemáticas.
%
% Esta plantilla ha sido desarrollada desde cero, pero utiliza algunas partes
% del código de la plantilla original utilizada en apuntesDGIIM
% (https://github.com/libreim/apuntesDGIIM), basada a su vez en las plantillas
% 'Short Sectioned Assignment' de Frits Wenneker (http://www.howtotex.com),
% 'Plantilla de Trabajo' de Mario Román y 'Plantilla básica de Latex en Español'
% de Andrés Herrera Poyatos (https://github.com/andreshp). También recoge
% ideas de la plantilla 'Multi-Purpose Large Font Title Page' de
% Frits Wenneker y Vel (vel@latextemplates.com).
%
% Licencia:
% CC BY-NC-SA 4.0 (https://creativecommons.org/licenses/by-nc-sa/4.0/)
%
%%%%%%%%%%%%%%%%%%%%%%%%%%%%%%%%%%%%%%%%%%%%%%%%%%%%%%%%%%%%%%%%%%%%%%%%%


% ---------------------------------------------------------------------------
% CONFIGURACIÓN BÁSICA DEL DOCUMENTO
% ---------------------------------------------------------------------------

%\documentclass[11pt, a4paper, twoside]{article} % Usar para imprimir
\documentclass[10pt, a4paper]{article}

\linespread{1.3}            % Espaciado entre líneas.
\setlength\parindent{0pt}   % No indentar el texto por defecto.
\setlength\parskip{7pt}

% ---------------------------------------------------------------------------
% CONFIGURACIÓN PERSONALIZADA
% ---------------------------------------------------------------------------

\usepackage{config}

% ---------------------------------------------------------------------------
% PAQUETES BÁSICOS
% ---------------------------------------------------------------------------

% IDIOMA
\usepackage[spanish, es-tabla, es-lcroman, es-noquoting, es-minimal]{babel}

% FUENTES
\usepackage{newpxtext, newpxmath}   % Fuente similar a Palatino
\usepackage{FiraSans}                 % Fuente sans serif
\usepackage[T1]{fontenc}
\usepackage[italic]{mathastext}     % Utiliza la fuente del documento
                                    % en los entornos matemáticos

% MÁRGENES
\usepackage[margin=2.5cm, top=3cm]{geometry}

%  OTROS
\usepackage{hyperref}   % Enlaces
\usepackage{graphicx}   % Permite incluir gráficos en el documento

% ---------------------------------------------------------------------------
% FUENTE DE MATEMÁTICAS ADAPTABLE EN SANS-SERIF
% ---------------------------------------------------------------------------

\usepackage{sansmath}

\newif\IfInSansMode
\let\oldsf\sffamily
\renewcommand*{\sffamily}{\oldsf\sansmath\InSansModetrue}

% ---------------------------------------------------------------------------
% DISEÑO DE PÁGINA
% ---------------------------------------------------------------------------

\usepackage{pagecolor}
\usepackage{afterpage}

% ---------------------------------------------------------------------------
% CABECERA Y PIE DE PÁGINA
% ---------------------------------------------------------------------------

\usepackage{fancyhdr}   % Paquete para cabeceras y pies de página

% Indica que las páginas usarán la configuración de fancyhdr
\pagestyle{fancy}
\fancyhf{}

% Representa la sección de la cabecera
\renewcommand{\sectionmark}[1]{\markboth{\thesection\quad #1}{}}

% Parte derecha de la cabecera
\fancyhead[R]{\sffamily \thepage}

% Parte izquierda de la cabecera
\fancyhead[L]{\sffamily{\leftmark}}

% Elimina la línea de la cabecera
\renewcommand{\headrulewidth}{0pt}

% Controla la altura de la cabecera para que no haya errores
\setlength{\headheight}{14pt}

% ---------------------------------------------------------------------------
% TÍTULOS DE PARTES Y SECCIONES
% ---------------------------------------------------------------------------

\usepackage{titlesec}

% Estilo de los títulos de las partes
\titleformat{\part}[hang]{\Huge\bfseries\sffamily}{\thepart\hspace{20pt}\textcolor{500}{|}\hspace{20pt}}{0pt}{\Huge\bfseries}
\titlespacing*{\part}{0cm}{-2em}{2em}[0pt]

% Reiniciamos el contador de secciones entre partes (opcional)
\makeatletter
\@addtoreset{section}{part}
\makeatother

% Estilo de los títulos de las secciones, subsecciones y subsubsecciones
\titleformat{\section}
  {\Large\bfseries\sffamily}{\thesection}{1em}{}

\titleformat{\subsection}
  {\Large\sffamily}{\thesubsection}{1em}{}[\vspace{.5em}]

\titleformat{\subsubsection}
  {\sffamily}{\thesubsubsection}{1em}{}

% ---------------------------------------------------------------------------
% ENTORNOS PERSONALIZADOS
% ---------------------------------------------------------------------------

\usepackage{mdframed}

%% MARCOS DE LOS ESTILOS

% Configuración general de mdframe, los estilos de los teoremas, etc
\mdfsetup{
  skipabove=1em,
  skipbelow=1em,
  innertopmargin=1em,
  innerbottommargin=1em,
  splittopskip=2\topsep,
}

% Definimos los marcos de los estilos

\mdfdefinestyle{nth-frame}{
	linewidth=2pt, %
	linecolor= 500, %
	topline=false, %
	bottomline=false, %
	rightline=false,%
	leftmargin=0em, %
	innerleftmargin=1em, %
  innerrightmargin=1em,
	rightmargin=0em, %
}%

\mdfdefinestyle{nprop-frame}{
	linewidth=2pt, %
	linecolor= 300, %
	topline=false, %
	bottomline=false, %
	rightline=false,%
	leftmargin=0pt, %
	innerleftmargin=1em, %
	innerrightmargin=1em,
	rightmargin=0pt, %
}%

\mdfdefinestyle{ndef-frame}{
	linewidth=2pt, %
	linecolor= 500, %
	backgroundcolor= 50,
	topline=false, %
	bottomline=false, %
	rightline=false,%
	leftmargin=0pt, %
	innerleftmargin=1em, %
	innerrightmargin=1em,
	rightmargin=0pt, %
}%

\mdfdefinestyle{ejer-frame}{
	linewidth=2pt, %
	linecolor= 300, %
	backgroundcolor= 50,
	topline=false, %
	bottomline=false, %
	rightline=false,%
	leftmargin=0pt, %
	innerleftmargin=1em, %
	innerrightmargin=1em,
	rightmargin=0pt, %
}%

\mdfdefinestyle{ejemplo-frame}{
	linewidth=0pt, %
	linecolor= 300, %
	leftline=false, %
	rightline=false, %
	leftmargin=0pt, %
	innerleftmargin=1.3em, %
	innerrightmargin=1em,
	rightmargin=0pt, %
	innertopmargin=0em,%
	innerbottommargin=0em, %
	splittopskip=\topskip, %
}%

% Asignamos los marcos a los estilos
\surroundwithmdframed[style=nth-frame]{nth}
\surroundwithmdframed[style=nprop-frame]{nprop}
\surroundwithmdframed[style=nprop-frame]{ncor}
\surroundwithmdframed[style=ndef-frame]{ndef}
\surroundwithmdframed[style=ejer-frame]{ejer}
\surroundwithmdframed[style=ejemplo-frame]{ejemplo}
\surroundwithmdframed[style=ejemplo-frame]{sol}

%%%%%%%%%%%%%%%%%%%%%%%%%%%%%%%%%%%%%%%%%%%%%%%%%%%%%%%%%%%%%%%%%%%%%%%%%%%%%
% ---------------------------------------------------------------------------
% COMIENZO DEL DOCUMENTO
% ---------------------------------------------------------------------------
%%%%%%%%%%%%%%%%%%%%%%%%%%%%%%%%%%%%%%%%%%%%%%%%%%%%%%%%%%%%%%%%%%%%%%%%%%%%%

\begin{document}

% ---------------------------------------------------------------------------
% PORTADA EXTERIOR
% ---------------------------------------------------------------------------

\newpagecolor{500}\afterpage{\restorepagecolor} % Color de la página
\begin{titlepage}

  % Título del documento
	\parbox[t]{\textwidth}{
			\raggedright % Texto alineado a la izquierda
			\fontsize{50pt}{50pt}\selectfont\sffamily\color{white}{
			  \textbf{\asignatura}
      }
	}

	\vfill

	%% Autor e información del documento
	\parbox[t]{\textwidth}{
		\raggedright % Texto alineado a la izquierda
		\sffamily\large\color{white}
		{\Large \autor }\\[4pt]
		\grado\\
		\universidad\\[4pt]
		\texttt{\enlaceweb}
	}

\end{titlepage}

% ---------------------------------------------------------------------------
% PÁGINA DE LICENCIA
% ---------------------------------------------------------------------------

\thispagestyle{empty}
\null
\vfill

%% Información sobre la licencia
\parbox[t]{\textwidth}{
  \includegraphics{by-nc-sa.pdf}\\[4pt]
  \raggedright % Texto alineado a la izquierda
  \sffamily\large
  {\Large Este libro se distribuye bajo una licencia CC BY-NC-SA 4.0.}\\[4pt]
  Eres libre de distribuir y adaptar el material siempre que reconozcas a los\\
  autores originales del documento, no lo utilices para fines comerciales\\
  y lo distribuyas bajo la misma licencia.\\[4pt]
  \texttt{creativecommons.org/licenses/by-nc-sa/4.0/}
}

% ---------------------------------------------------------------------------
% PORTADA INTERIOR
% ---------------------------------------------------------------------------

\begin{titlepage}

  % Título del documento
	\parbox[t]{\textwidth}{
			\raggedright % Texto alineado a la izquierda
			\fontsize{50pt}{50pt}\selectfont\sffamily\color{500}{
			  \textbf{\asignatura}
      }
	}

	\vfill

	%% Autor e información del documento
	\parbox[t]{\textwidth}{
		\raggedright % Texto alineado a la izquierda
		\sffamily\large
		{\Large \autor}\\[4pt]
		\grado\\
		\universidad\\[4pt]
		\texttt{\enlaceweb}
	}

\end{titlepage}

% ---------------------------------------------------------------------------
% ÍNDICE
% ---------------------------------------------------------------------------

\thispagestyle{empty}
\tableofcontents
\newpage

% ---------------------------------------------------------------------------
% CONTENIDO
% ---------------------------------------------------------------------------

%\part{Teoría}
%\section{Ecuaciones lineales: teoremas de existencia y unicidad.}


\begin{ndef}[Ecuación diferencial lineal]
Sea $(\alpha,\beta)\subseteq \R$ un intervalo abierto y sean $A:(\alpha,\beta)\to \M_d(\R)$ y $b:(\alpha,\beta)\to \R^d$ funciones continuas. Entonces una ecuación diferencial lineal es de la forma:

\begin{equation}
x'=A(t)x+b(t) \tag{C} \label{lin:completa}
,\end{equation}

y se dice \emph{completa}, o bien

\begin{equation}
x'=A(t)x \tag{H} \label{lin:homogenea}
,\end{equation}

y en este caso se dice \emph{homogénea}.
\end{ndef}

\begin{nth}[Teorema de existencia y unicidad de la solución]
Dados $t_0\in(\alpha,\beta)$ y $x_0 \in \R^d$ y consideramos el \emph{problema de valores intermedios} (PVI):

\begin{equation}
\left\{\begin{array}{rl}
x' &=A(t)x+b(t), \\
x(t_0)&=x_0.
\end{array}\right.
\tag{PVI}\label{lin:PVI}
\end{equation}

Entonces, existe una única solución $\varphi : (\alpha,\beta)\to\R^d$, con $\varphi \in \C ^ 1(\R)$, que verifica

\[
\varphi'(t) = A(t)\varphi(t) + b(t) \qquad \forall t \in (\alpha,\beta)
\]

y que además cumple $\varphi(t_0)=x_0$, es decir, una única solución de \ref{lin:PVI}.
\end{nth}

Aplicando el teorema al PVI

\[
\begin{cases}
  x' = Ax, \\
  x(t_0) = e_j.
\end{cases}
\tag{PVI$_j$} \label{lin:PVIj}
\]

para cada $j=1,\dots,d$, obtenemos soluciones $\varphi_j$. Sea $x_0\in \R^d$, y consideramos $\varphi = \sum_{j=1}^d \langle x_0,e_j \rangle\varphi_j$.
Es inmediato comprobar que $\varphi$ es solución de \eqref{lin:PVI} con $b=0$. Por tanto, hemos probado:

\begin{ncor}
  Las soluciones de \eqref{lin:homogenea} son un espacio vectorial de dimensión $d$.
\end{ncor}

Ahora, consideremos $\varphi$ la solución de \eqref{lin:PVI}, para $x_0\in\R^d$ arbitrario, es decir, una solución cualquiera de \eqref{lin:completa}.
Sean $\varphi_c$ la solución de

\begin{equation}
\left\{\begin{array}{rl}
x' &=A(t)x+b(t), \\
x(t_0)&=x_0',
\end{array}\right.
\end{equation}

con $x_0'\in\R^d$ arbitrario, y $\varphi_h$ la de

\begin{equation}
\left\{\begin{array}{rl}
x' &=A(t)x, \\
x(t_0)&=x_0 - x_0',
\end{array}\right.
\end{equation}

y ahora es rutinario comprobar que $\varphi = \varphi_c + \varphi_h$. Además, para cualquier otra solución $\psi_{h}$ de \eqref{lin:homogenea}, $\varphi_c + \psi_h$
es solución de \eqref{lin:completa}.
En esta ocasión, hemos probado:

\begin{ncor}
  El conjunto $S_c$ de soluciones de \eqref{lin:completa}, dada cualquier solución $\varphi$ de la misma, es el espacio afín

  \[
    S_c = \varphi + S_h,
    \]

    con $S_h$ el conjunto de soluciones de \eqref{lin:homogenea}.
\end{ncor}

\subsection{Matriz Fundamental Principal en un punto}
\begin{ndef}
Dada una Ecuación Diferencial Ordinaria (EDO) lineal y homogénea con coeficientes constantes:
\[
\begin{cases}
x'=A(t)x\\
x(t_0)=x_0
\end{cases}
\]
Definimos la matriz fundamental principal(m.f.p) en $t_0$ como: $$ \phi(t)=e^{A(t-t_0)} \qquad \forall t \in (\alpha,\beta)$$
\end{ndef}

Ahora nos preguntamos cómo se calcula la m.f.p. cuando la matriz $A$ es diagonalizable. En este caso, $\exists P, D \in \M_d(\C)$ tales que $det(P)\neq 0, \quad D=diag(\lambda_1,...\lambda_d)$ y $A=P\cdot D \cdot P^{-1}$. En cuyo caso, la matriz fundamental principal será:
$$e^{At}=P\cdot e^{Dt}\cdot P^{-1}=P \cdot  \begin{pmatrix}e^{\lambda_1 t} & \cdots & 0 \\ \vdots & \ddots & \vdots \\0 & \cdots & e^{\lambda_d t}\end{pmatrix} \cdot P^{-1}$$


Cuando $A$ no sea diagonalizable tendremos que utilizar la forma canónica de Jordan. En este caso, $\exists P,J\in \M_d(\C)$ tales que $det(P) \neq 0 $ y $J$ diagonal por bloques:
$$\begin{pmatrix}J_1 & & \\
 & & \\
 & & J_r \end{pmatrix}$$
 donde $orden(J_k)=1 \quad \forall =1,...,r \implies J_k=(\lambda_k)$ con $\lambda_k$ valor propio de $A$ o si $orden(J_k)>1$,
 $$J_k=
 \begin{pmatrix}
 \lambda_k & 1  &\cdots & 0 \\
  0 & \lambda_k & \ddots & \vdots \\
  \vdots & \ddots & \ddots & 1 \\
  0 & \cdots & 0 & \lambda_k
\end{pmatrix}$$

FIXME : POR TERMINAR
\subsection{Estabilidad en el sentido de Lyapunov}
%% CLASE LUNES 17

\begin{nprop}
Las siguientes afirmaciones son equivalentes:
\begin{nlist}
\item Todas las soluciones de \eqref{lin:completa} son atractores.
\item Existe una solución de \eqref{lin:completa} que es un atractor.
\item La solución trivial $y = 0$ de \eqref{lin:homogenea} es un atractor.
\item Todas las soluciones de \eqref{lin:homogenea} convergen hacia el vector 0 cuando $t \to +\infty$.
\item La matriz fundamental de \eqref{lin:homogenea} principal en $t_0$ converge hacia la matriz 0 cuando $t \to +\infty$.
\end{nlist}
\end{nprop}

\begin{ncor}
Los atractores de la ecuación \eqref{lin:homogenea} son asintóticamente estables.
\end{ncor}

\begin{ndef}
\begin{itemize}
\item Se dice que la ecuación \eqref{lin:completa} es estable si todas su soluciones son estables
\item Se dice que la ecuación \eqref{lin:completa} es asintóticamente estable si todas sus soluciones son asintóticamente estables.
\end{itemize}
\end{ndef}

No siempre podemos calcular la matriz fundamental, pero nosotros nos centraremos en el caso escalar donde sí podemos calcularla.

\subsection{Estabilidad de ecuaciones lineales escalares}
Sean $a,b:(\alpha,+\infty)\to \R$ continuas. Se considera la ecuación diferencial lineal escalar \eqref{completa} y sea $t_0\in(\alpha,+\infty)$.
La matriz fundamental principal en $t_0$ en $\phi(t)=exp(\int^t_{t_0} a(s)ds)$ y por tanto podemos caracterizar la estabilidad de \eqref{completa} controlando una primitiva del coeficiente $a(t)$.
\begin{nth}{Proposicion 3}
\begin{itemize}
\item La ecuación de \ eqref{completa} es estable sii la función $a(t)$ tiene una primitiva acotada superiormente en $[t_0,+\infty]$.
\item La ecuación \ eqref{completa} es a.e. sii la función $a(t)$ tiene una primitiva que converge hacia $-\infty$ [....]
\end{itemize}
\end{nth}
Ejemplo 1:
- Inestable
- A.E.
- estable
-
- Inestable
- estable
- estable
-


Estable $-> \lambda <= 0$
A.E. $-> \lambda < 0$
Inestable $-> \lambda > 0$

\subsection{Estabilidad de ecuaciones lineales con coeficientes constantes}
Sea $A \in \M_d(\R)$ una matriz cuadrada.  Usaremos la siguiente notación:
NOTACIÓN:
\begin{ndef}{Espectro de $A$}
El espectro de $A$ es el conjunto de valores propios de $A$, tanto reales como complejos, $\sigma(A)={\lambda_1,...,\lambda_d}$, contados con su multiplicidad en nuestro caso.
\end{ndef}
\begin{ndef}{Multiplicidad del valor propio $\lambda_j$}
$m(\lambda_j) \forall \lambda_j \in \sigma{A}$
\end{ndef}
\begin{ndef}{Dimensión de cada subespacio propio $E_{\lambda_j}$}
$dim E_{\lambda_j} = dim ker(A-\lambda_{jI})= d-rango(A-\lambda I)$
\end{ndef}
\begin{ndef}
Los valores propios de A cuya parte real es 0 son:
$\sigma_0(A)={\lambda \in \sigma(A) \ Re(\lambda)=0}$
\end{ndef}

Consideramos la EDO lineal homogénea y autónoma $x'=Ax x\in\R^d$. El principal indicador para determinar la estabilidad de * es el máximo de las partes reales de los valores propios de A:
$\mu(A)=máx{Re(\lambda : \lambda\in\sigma(A))}$

\begin{nth}
\begin{enumerate}
\item Si $\mu(A)<0$ entonces la EDO lineal * es A.E.
\item Si $\mu(A)=0$ y $m(\lambda)=dimE_\lambda \forall \lambda \in \sigma_0(A)$, entonces la EDO lineal * es estable (pero no A.E.)
\item Si $\mu(A)=0$ y $\exists \lambda \in \sigma_0(A)$ tal que $m(\lambda) != dim E_lambda$, entonces la EDO lineal * es inestable.
\item Si $\mu(A)>0$ entonces la EDO lineal * es inestable.
\end{enumerate}
\end{nth}

A continuación, haremos una serie de ejemplos y proseguiremos con un poquito de teoría para poder producir la demostración de este enunciado.
Ejemplo 1:
{
\begin{equation}
\left\{\begin{array}{lcl}
x_1'=2x_1+3x_2\\
x_2'=5x_1-x_2
\end{array}\right.
\end{equation}
()=()()
\begin{equation}
\left\{\begin{array}{lcl}
x_1' = -2 3   x_1'
x_2'= 5 -1    x_2'
\end{array}\right.
\end{equation}

EC. $\lambda^2-traza(A)\lambda + det(A=0)$
$\lambda+3\lambda-13=0$
$\sigma(A)={\frac{-3+\sqrt{61}}{2},\frac{-3-\sqrt{61}}{2}}={2.4...,5.4...}$
$\mu(A)>2.4 >0$
La ecuación es inestable.


%% MARTES 18

Recordemos la forma canónica de Jordan.

Sea $A \in \M_d(\R) \exists P, J \in \M_d(\C) tales que:$
\begin{itemize}
\item $P$ es regular (invertible)
\item $J$ es diagonal por bloques:
\begin{equation}
J=diag(J_1,...,J_k) J_j matriz diagonal \forall j=1,...,k
\end{equation}
\item $A=P\cdot J\cdot P^{-1}$
Para calcular $J$ se usan:
\item $\sigma(A)={\lambda_1,...\lambda_d}$
\item $m(\lambda_j)$
\item $dim E_{\lambda_j}$
\item $\nu(\lambda_j)=min{k\in\mathbb{Z}^+: BORRA}$
Los bloques de $J$ pueden ser:
\item de orden 1: $J_k=(\lambda_k) \lambda_k \in \sigma(A)$
\item de orden mayor que 1: $J_k=diag(\lambda_k,...\lambda_k) (1...1 encima de la diag) con \lambda_k \in \sigma(A)$
\item el orden $(J_k) \leq \nu(\lambda_k)$
\item la suma de los órdenes de todos los bloques asociados aun mismo valor propio es igual a la multiplicidad de ese valor propio
\end{itemize}

INSERTE EJEMPLO DE CAJA DE JORDAN

Con estos datos, podemos hacer el siguiente cálculo de la exponencial:

$$A=P\cdot J \cdot P^{-1} \implies A^n=P \cdot J^n \cdot P^{-1} \implies e^A=P \cdot e^J \cdot P^{-1}$$

Y de manera análoga:
\begin{itemize}
\item $e^{At}=P \cdot e^{Jt} \cdot P^{-1}$
\item $e^{Jt}=diag(e^{J_1t},...,e^{J_rt} )$
\end{itemize}

%% Inserte ejemplo

\begin{lema}
$$|||e^{Jt}|||_1=max{|||e^{J_1t}|||_1,|||e^J_rt|||_1}$$
\end{lema}

\begin{lema}
$\exists T \geq 0 $ tal que :
$$|||e^{Jt}|||_1=(1+...+\frac{t^n}{n!})\cdot e^{\mu t} \qquad \forall t \geq T$$
donde:
$$\mu=max{Re(\lambda_j):\lambda_j\in \sigma(A)}$$
$$n=max{\nu(\lambda_j)-1 : \lambda_j\in \sigma(A) y Re(\lambda_j)=\mu}$$
\end{lema}

Usando el lema anterior:
%%\begin{equation}
%%|||e^{Jt}|||_1=max
%%\left\{\(\begin{array}{lcl}
%%1+...+\frac{t^{\nu(\lambda_j)-1}}{(\nu (\lambda_j)-1)!}
%%\end{array}\right\)
%%\cdot e^{Re(\lambda_j)t} : \lambda_j \in \sigma(A) \right\}
%%\tag{cosarara}\label{cosarara}
%%\end{equation}
%%$\exists T \geq 0$ tal que:
%%\begin{equation}
%%(1+...+\frac{t^{\nu(\lambda_j)-1}}{\nu(\lambda_j - 1})\cdot e^{Re(\lambda_j)t} \leq (1+...1+\frac{t^n}{n!}e^{\mu t})
%%\end{equation}

Finalmente, podemos demostrar el teorema anterior:

\begin{proof}
$\exists P, J \in \M_d(\C)$ tales que $P$ es inestable, $J$ es diagonal por bloques y $A=P\cdot J \cdot P^{-1}$. Por tanto $$e^{At}=P\cdot e^{Jt} \cdot P^{-1}$$.
Defino en $\M_d(\C)$ la norma matricial:
$$|||B|||=|||P\cdot J \cdot P^{-1}|||_1$$.
Así, tenemos que:
\begin{enumerate}
\item Si $\mu < 0 \implies \lim_{t\to\infty}|||e^{At}|||=0\implies x'=Ax$ es A.E.
\item Si $\mu = 0 y dimE_j=m(\lambda)' \forall \lambda \in \sigma_0(A)\implies n=0 \implies |||e^{At}|||=1\implies|||e^{At}||| es estable pero no A.E.$
\item Si $\mu = 0 y \exists \lambda in \sigma_0(A) \nu(A)>1 \implies n\geq 1\implies |||e^{At}|||\geq 1+t$ no es acotada en $[0,+\infty)$
\item Si $\mu > 0 \qquad |||e^{At}|||\geq e^{\mu t}\to \infty$ no es acotada en $[0,+\infty)$
\end{enumerate}
\end{proof}


%% MIÉRCOLES 19
\section{Órbitas}
%% Lunes 24
\subsection{Diagramas de fase planos}
Consideramos la EDO lineal homogénea autónoma en $\R^2$:
\begin{equation}
\left\{\begin{array}{rl}
x_1'=a_{11}x_1+a_{12}x_2 \\
x_2'=a_{21}x_1+a_{22}x_2
\end{array}\right.\Leftrightarrow
\begin{pmatrix}
x_1 \\
x_2 \\
\end{pmatrix}
'=
\begin{pmatrix}
a_{11} & a_{12} \\
a_{21} & a_{22} \\
\end{pmatrix}
\begin{pmatrix}
x_1 \\
x_2 \\
\end{pmatrix}
\end{equation}

\begin{itemize}
\item Caso 1: Cuando los valores propios de la matriz A son reales: $\sigma(A)={\lambda_1,\lambda_2}\subset \R$
\begin{enumerate}
\item Son reales, positivos y distintos: $0<\lambda_1<\lambda_2 \Leftrightarrow \left\{ \begin{array}{ll}\Delta > 0 \\ traza(A)>0 \\ det(A) > 0 \end{array}\right.$
Inserte dibujito de fuente (flechas hacia fuera)
$x(t)=c_1 e^{\lambda_1t}v_1+c_2 e^{\lambda_2t}v_2 \qquad v_j \in E_{\lambda_j}$
\item Sonreales, negativos y distintos: $\lambda_2 < \lambda_1 \Leftrightarrow \left\{ \begin{array}{ll}\Delta > 0 \\ traza(A) < 0 \\ det(A)>0\end{array}\right.$
Inserte dibujito de sumidero (flechas hacia dentro)
\item Son reales, no nulos y con distinto signo : $\lambda_1 < 0 < \lambda_2 \Leftrightarrow det(A)<0$
Inserte dibujito de punto hiperbólico o punto silla
\item Son reales, positivos e iguales : $\lambda_1=\lambda_2>0 \Leftrightarrow \left\{ \begin{array}{ll} \Delta = 0 \\ traza(A)>0\end{array}\right.$.
\begin{enumerate}
\item Si $A$ es diagonalizable : En este caso tendremos todo el espacio y volveremos a tener una fuente.
\item Si $A$ no es diagonalizable : Nos encontramos con una fuente degenerada.
\end{enumerate}
\item Son reales, negativos e iguales : $\lambda_1=\lambda_2<0 \Leftrightarrow \left\{ \begin{array}{ll} \Delta = 0 \\ traza(A)<0\end{array}\right.$.
\begin{enumerate}
\item Si $A$ es diagonalizable : En este caso tendremos todo el espacio y volveremos a tener un sumidero.
\item Si $A$ no es diagonalizable : Nos encontramos con un sumidero degenerado.
\end{enumerate}
\item Un valor propio es 0 y el otro positivo $\lambda>0$
Inserte dibujito rectas paralelas y perpendiculares a una única recta generada por el 0(Hacia dentro). Rectas  de puntos de equilibrio estables
\item Un valor propio es 0 y el otro negativo $\lambda<0 $
Inserte dibujito rectas paralelas y perpendiculares a una única recta generada por el 0(Hacia fuera). Rectas  de puntos de equilibrio inestables
\item Dos valores propios iguales a 0 $\lambda_1=\lambda_2=0 \Leftrightarrow \left\{ \begin{array}{ll}\Delta=0 \\ traza(A)=0 \\ det(A)=0\end{array}\right.$
\begin{enumerate}
\item Si $A$ es diagonalizable $A=0$
\item Si $A$ no es diagonalizable.
\end{enumerate}
\end{enumerate}
\item Caso 2: Cuando los valores propios de la matriz A no son reales : $\sigma(A)={\lambda_1,\lambda_2}\subset \mathbb{C} - \R \quad \lambda_1=a+bi \; \lambda_2=a-bi$
\begin{enumerate}
\item Tiene parte real positiva $\Delta < 0 \\ traza(A) > 0$ Fuente en espiral
\item Tiene parte real negativa $\Delta < 0 \\ traza(A) < 0$ Sumidero en espiral
\item Tiene parte real igual a 0. Diagrama de Poincaré
\end{enumerate}
\end{itemize}

%\newpage

\part{Ejercicios}
\section{Tema 1}

\begin{ejer}
  Definir la distribución de probabilidad, función de distribución y función característica de una variable aleatoria y expresar dichas funciones en términos de la función masa de probabilidad o la función de densidad, según que la variable sea de tipo discreto o continuo, respectivamente.
\end{ejer}

\begin{sol}
  La \emph{distribución de probabilidad} de una variable aleatoria es una función de probabilidad definida en el espacio de Borel:
  \begin{align*}
    P_X : \B & \to \R \\
    B & \mapsto P_X(B) = P(X^{-1}(B)) = P[X \in B]
  \end{align*}

  La \emph{función de distribución} es la función de puntos definida como:
  \begin{align*}
    F: \R & \to \R \\
    x & \mapsto  F_X(x) = P_X((-\infty, x]) = P[X \leq x]
  \end{align*}

  La \emph{función característica} es la función definida como $\varphi_X(t) = E[e^{itX}], \; \; \forall t \in \R$.

  \begin{itemize}
    \item Caso discreto:
    \begin{multline*}
      \forall x \in \R, \; F_X(x) = P_X[(-\infty, x]] = P_X[\{x_1, \ldots, x_n, \ldots \}] = P[X \in \{x_1,  \ldots, x_n, \ldots \}] = \\
      = P[\bigcup \limits^\infty_{n = 1} (X = x_n)] = \sum \limits^\infty_{i = 1} P[X = x_i] = \sum \limits^\infty_{i = 1} f(x_i)
    \end{multline*}

    $$\forall t \in \R, \; \varphi_X(t) = E[e^{itX}] = \sum \limits^\infty_{i = 1} e^{itx_i}P[X = x_i] = \sum \limits^\infty_{i = 1} e^{itx_i} f(x_i)$$

    \item Caso continuo:
    $$\forall x \in \R, F_X(x) = P_X[(-\infty,x]] = P[X \in (-\infty, x]] = P[X \leq x] = \int^x_{-\infty} f(y)dy$$

    $$\forall t \in \R, \; \varphi_X(t) = E[e^{itX}] = \int_\R e^{itx} f(x) dx$$
  \end{itemize}
  Donde $f(x)$ representa la función masa de probabilidad en el caso discreto, y la función densidad en el caso continuo.
\end{sol}

\begin{ejer}
  Demostrar que la distribución de probabilidad de una variable aleatoria es una medida de probabilidad definida sobre la $\sigma$-álgebra de Borel.
\end{ejer}

\begin{sol}
  Veamos si
  \begin{align*}
    P_X : \B & \to \R \\
    B & \mapsto P_X(B) = P(X^{-1}(B)) = P[X \in B]
  \end{align*}
  cumple con los tres Axiomas de Kolmogorov, usando que P es una probabilidad:

  \begin{nlist}
    \item $P_X(B) = P[X \in B] \geq 0$, $\forall B \in \B$.
    \item $\forall \{B_n\}_{n \in \N} \subset \B$ disjuntos $$ P_X(\bigcup \limits^\infty_{n = 1} B_n) = P[X \in \bigcup \limits^\infty_{n = 1} B_n] = P[\bigcup \limits^\infty_{n = 1} (X \in B_n)] = \sum \limits^\infty_{n = 1} P[X \in B_n] = \sum \limits^\infty_{n = 1} P_X(B_n)$$
    \item $P_X(\R) = P[X^{-1}(\R)] = P[\W] = 1$
  \end{nlist}
\end{sol}

\begin{ejer}
  Demostrar la caracterización de vectores aleatorios mediante variables aleatorias.
\end{ejer}

\begin{sol}
  Sea $X = (X_1, \ldots, X_n) : (\W, \A, P) \to (\R^n, \B^n)$, entonces $X$ es vector aleatorio $\iff$ $X_1, \ldots, X_n$ son variables aleatorias.

  \begin{multline*}
    X \text{ v.a} \iff \forall x = (x_1, \ldots, x_n) \in \R^n, \; X^{-1}((-\infty, x]) = [X \leq x] = \\ =  [X_1 \leq x_1, \ldots, X_n \leq x_n] = \bigcap \limits^n_{i = 1} [X_i \leq x_i] \subset \A \iff \forall i \in \{1, \ldots, n\}, \; \forall x_i \in \R \\ \A \supset [X_i \leq x_i] = X^{-1}((-\infty, x_i]) \iff X_1, \ldots, X_n \text{ v.a }
  \end{multline*}

  Vemos la penúltima implicación tomando $x_j = \infty \forall j \neq i$ en $\bigcap \limits^n_{i = 1} [X_i \leq x_i] \subset \A \implies [X_i \leq x_i] \subset \A$, y al revés sabiendo que $\bigcap \limits^n_{i = 1} [X_i \leq x_i] \subset [X_i \leq x_i] \subset \A$.
\end{sol}

\begin{ejer}
  Sean $X$ e $Y$ variables aleatorias independientes tales que $E[X] = 0$, $Var(X) = 1$ y la variable $Y$ tiene distribución uniforme en $[-\pi, \pi]$. Consideremos el proceso estocástico $\{X_t\}_{t \geq 0}$ definido por

  $$X_t = X \cos(t + Y)$$

  Calcular las funciones media y covarianza y decir si el proceso es débilmente estacionario.
\end{ejer}

\begin{sol}
  Primero veamos que efectivamente $\{X_t\}_{t \geq 0}$ es un proceso estocástico. Para $t \geq 0$, tenemos que $X_t : (\W, \A, P) \to (\R, \B)$, luego falta ver que es medible, pero funciones Borel de variables aleatorias son medibles y el producto de funciones medibles es medible.

  Función media: $\mu(t) = E[X_t] = E[X \cos(t + Y)]$, como $\cos(t + Y)$ es función medible por composición de funciones medibles junto a que $X$ y $Y$ son independientes, entonces $X$ y $\cos(t + Y)$ son independientes; luego $E[X \cos(t + Y)] = E[X] E[\cos(t + Y)] = 0$. Luego $\mu(t) = 0$, $\forall t \geq 0$.

  Función covarianza: $C(s,t) = R(s,t) - \mu_t \mu_s = R(s,t) = E[X_t X_s] = E[X^2 \cos(t + Y) \cos(s + Y)]$, por la independencia otra vez, $E[X^2 \cos(t + Y) \cos(s + Y)] = E[X^2] E[\cos(t + Y) \cos(s + Y)] = Var(X) E[\cos(t + Y) \cos(s + Y)] = E[\cos(t + Y) \cos(s + Y)]$, resolvemos la integral $\frac{1}{2\pi}\int^\pi_{-\pi} \cos(t + y) \cos(s + y) dy$ y tenemos que $C(s,t) = \frac{1}{2} \cos(s - t)$.


  Este pe. es débilmente estacionario ya que su función media es constante, $\mu(t) = 0$, $\forall t \geq 0$; y la función covarianza cumple $C(s,t) = \frac{1}{2} \cos(s - t) = \frac{1}{2} \cos(s + h - t - h) = C(s + h, t + h) = C(0, t - s)$, $\forall t,s \geq 0$.
\end{sol}

\begin{ejer}
  Ejemplos de procesos estocásticos atendiendo a la clasificación según el espacio de estados y espacio paramétrico.
\end{ejer}

\begin{sol}
  Ejemplos:

  \begin{itemize}
    \item PDTD: Cantidad de material en el almacén en el día n-ésimo.
    \item PCTD: Velocidad de un coche en el minuto n-ésimo.
    \item PDTC: Nº de coches en el instante $t$.
    \item PCTC: Posición de una particula en el instante $t$.
  \end{itemize}
\end{sol}

\begin{ejer}
  Demostrar que si las componentes del vector aleatorio $X = (X_1, \ldots, X_n)^T$ son independientes, entonces la funcion característica del vector $X$ es igual al producto de las funciones características de las variables $X_k$, $k = 1, \ldots, n$.
\end{ejer}

\begin{sol}
  Sea un vector aleatorio $X = (X_1, \ldots, X_n)$, la función característica del vector $X$ es $\varphi_X(t) = E[e^{it^TX}] = E[e^{i(t_1 X_1 + \ldots + t_n X_n)}] = E[e^{i t_1 X_1 + \ldots + i t_n X_n}] = E[\prod \limits^n_{k = 1} e^{i t_k x_k}] = \prod \limits^n_{k = 1} E[e^{i t_k x_k}] = \prod \limits^n_{k = 1} \varphi_{X_k} (t_k)$. Podemos separar las esperanzas ya que $X_1, \ldots, X_n$ son independientes y $e^{itx}$ es función Borel de variable aleatoria; y por tanto $e^{itX_k}, \; k = 1, \ldots, n$, siguen siendo independientes.
\end{sol}

\begin{ejer}
  Demostrar que para un proceso con incrementos independientes las distribuciones finito dimensionales del proceso (esto es, las distribuciones de los vectores ($X_{t_1}, X_{t_2}, \ldots, X_{t_n})^T$, $\forall t_1 < t_2 < \ldots < t_n$) están determinadas por las distribuciones marginales unidimensionales ($X_{t_1}$) y por las distribuciones de los incrementos ($X_{t_i} - X_{t_1 - 1}$, $i = 2, \ldots, n$).
\end{ejer}

\begin{sol}
  Tenemos un $\{X_t\}_{t \in T}$ p.e con incrementos independientes, entonces $\forall t_1 < \ldots < t_n \in T$, $X_{t_1}, X_{t_2} - x_{t_1}, \ldots, X_{t_n} - X_{t_{n-1}}$ son v.a independientes.

  Definimos $S_1 = X_{t_1}$, $S_2 = X_{t_2} - X_{t_1} \implies X_{t_2} = S_1 + S_2$

  $S_3 = X_{t_3} - X{t_2} \implies X_{t_3} = S_1 + S_2 + S_3$, por tanto $X_{t_k} = \sum \limits^k_{j = 1} S_j$.

  $\phi_{(X_{t_1}, \ldots, X_{t_n})}(u_1, \ldots, u_n) = \phi_{(s_1, \ldots, s_n)}(v_1, \ldots, v_n)$

\end{sol}

\begin{ejer}
  Demostrar que la clase de rectángulos medibles $\mathcal{C}^\N$ es una semi-álgebra.
\end{ejer}

\begin{sol}
  Sean los espacios medibles $(\W_n, \A_n), \; \forall n \in \N$, definimos $$\mathcal{C}^\N = \{\prod \limits^\infty_{n = 1} A_n : A_n \in \A_n \}$$.

  Veamos que $\mathcal{C}^\N$ es una semi-álgebra.
  \begin{nlist}
    \item $\W, \emptyset \in \mathcal{C}^\N$.

    Como $\emptyset, \W_i \in \A_n \; \forall n \in \N$ ($\A_n$ es $\sigma$-álgebra) $\implies \emptyset, \W \in \mathcal{C}^\N$.

    \item $\forall A, B \in \mathcal{C}^\N \implies A \cap B \in \mathcal{C}^\N$.

    Sean $A, B \in \mathcal{C}^\N$, $A \cap B = (\prod \limits^\infty_{n = 1} A_n) \cap (\prod \limits^\infty_{n = 1} B_n) = \prod \limits^\infty_{n = 1} (A_n \cap B_n) \in \mathcal{C}^\N$, ya que $A_i, B_i \in \A_i \implies A_i \cap B_i \in \A_i$ (por ser $\A_i$ $\sigma$-álgebra).

    \item $\forall A, B \in \mathcal{C}^\N$, $\exists \{C_i\}^n_{i = 1} \subset \mathcal{C}^\N$ disjuntos dos a dos tal que $A \setminus B = \bigcup \limits^n_{i = 1} C_i$.

    Sean $A, B \in \mathcal{C}^\N$, $A \setminus B = (\prod \limits^\infty_{n = 1} A_n) \setminus(\prod \limits^\infty_{n = 1} B_n) = \prod \limits^\infty_{n = 1} (A_n \setminus B_n) = C \in \mathcal{C}^\N$, ya que $\emptyset \in \A_i$ y $A_i \setminus B_i = A_i \cap \overline{B_i} \in A_i$ ($\A_i$ es $\sigma$-álgebra).
  \end{nlist}
\end{sol}

\begin{ejer}
  Demostrar el Teorema de medibilidad para procesos estocásticos en tiempo discreto.
\end{ejer}

\begin{sol}
  Sea la función \begin{align*}
    \chi : (\W, \A) & \to (\R^\N, \B^\N) \\
    \w & \mapsto X(\w) = \{X_n(\w)\}_{n \in \N}
  \end{align*}

  Entonces $\chi$ es medible $\iff \{X_n\}_{n \in \N}$ son medibles.
\end{sol}


\end{document}
